


%===========================================
% ESTE ARQUIVO É UMA INSTRUÇÃO/PROMPT DETALHADO PARA INTELIGÊNCIA ARTIFICIAL.
%
% OBJETIVO: Criar um template LaTeX COMPILÁVEL, TOTALMENTE MODULAR e livre de "gambiarras",
% utilizando EXCLUSIVAMENTE a classe MEMOIR para atender aos padrões ABNT e aos
% requisitos de numeração e estilo do usuário.
%
% OBSERVAÇÃO CRÍTICA: A modularização (criação de macros, \makechapterstyle, etc.)
% DEVE ser feita diretamente neste arquivo .tex, NÃO em um arquivo de classe (.cls)
% separado, para facilitar os testes e a depuração.
%===========================================
\documentclass[12pt,a4paper,oneside]{memoir}
\usepackage[brazilian]{babel}
\usepackage{fontspec}
\usepackage{tikz} % NECESSÁRIO PARA O ESTILO DE CAPÍTULO PERSONALIZADO
%===========================================

\setmainfont{CMU Serif}
\setlrmarginsandblock{3cm}{2cm}{*}
\setulmarginsandblock{3cm}{2cm}{*}
\checkandfixthelayout
\OnehalfSpacing % Espaçamento entre linhas de 1,5
\setlength{\parindent}{1.25cm} % Recuo do paragrafo
%===========================================


%===========================================
% INSTRUÇÃO IA: CONFIGURAÇÃO DE ESTILO DE CAPÍTULO MODULAR (O SWITCH)
%
% OBJETIVO: Implementar dois estilos de capítulo baseados na classe MEMOIR para
% aplicação modular e criar a macro de alternância.
%
% IMPLEMENTAÇÃO:
% 1. Estilo Numerado Principal (\makechapterstyle{CapitulosNumerado}): Implementar o código TikZ fornecido. (Aplicado aos capítulos numerados).
% 2. Estilo Não Numerado Simplificado (\makechapterstyle{estiloNaoNumerado}):
%    - A IA DEVE CRIAR ESTE ESTILO. Usado para capítulos \chapter* (Introdução, Referências, Listas, Apêndices/Anexos).
%    - Características: Título CENTRALIZADO, COM LINHA HORIZONTAL (como na ABNT), SEM O NÚMERO DECORATIVO GRANDE do TIkZ.
% 3. Comando Modular (\capituloCentralizado):
%    - Esta macro DEVE ALTERNAR temporariamente o estilo para 'estiloNaoNumerado', usar \chapter*{Título} (com as adições de TOC/numeração invisível), e retornar para 'CapitulosNumerado'.
%
% CRÍTICO: Garantir que o estilo simplificado obedeça a centralização e não cause conflitos.
%===========================================
\makechapterstyle{CapitulosNumerado}{
    \setlength{\beforechapskip}{0.46cm}
    \renewcommand{\printchaptername}{}
    \renewcommand{\chapternamenum}{}
    \renewcommand{\afterchapternum}{\par\nobreak\vskip 1.3cm}
    \renewcommand{\printchapternum}{%
        \noindent
        \begin{tikzpicture}[remember picture, overlay]
            \def\numsize{61}
            \def\barwidth{0.8pt}
            \def\leftoffset{2.3cm}

            \node[inner sep=-1pt, anchor=south west, opacity=0] (num) at (\leftoffset, 0)
             {\fontsize{\numsize}{\numsize}\selectfont\rmfamily\thechapter};
            \node[inner sep=0pt, anchor=south west] at (\leftoffset, -1pt)
             {\fontsize{\numsize}{\numsize}\selectfont\rmfamily\thechapter};
            \node[inner sep=0pt, anchor=south west, font=\Large\rmfamily] (chaptext)
                 at (0pt, -4pt) {\chaptername};

            \coordinate (StartPoint) at (0.5*\barwidth, 13pt);
            \draw[line width=\barwidth] (StartPoint) -- (StartPoint |- num.north) -- (num.north west);
            \coordinate (RightWall) at (\linewidth, 0);
            \draw[line width=\barwidth] (num.north east) -- (RightWall |- num.north) -- (RightWall |- num.south) -- (num.south east);
        \end{tikzpicture}%
    }
    \renewcommand{\printchaptertitle}[1]{%
        \raggedright\huge\bfseries ##1
    }
}
% A IA DEVE INSERIR AQUI A DEFINIÇÃO DE \makechapterstyle{estiloNaoNumerado} (simplificado e centralizado com a reta).

%===========================================
\newcommand{\folha}{
    \thispagestyle{empty}
    \vspace*{8cm}
    \begin{center}
        {\bfseries\scshape\Huge Folha de Rosto \par}
    \end{center}
    \clearpage
}
%===========================================
\newcommand{\contra}{
    \thispagestyle{empty}
    \vspace*{8cm}
    \begin{center}
        {\bfseries\scshape\Huge Contra Capa \par}
    \end{center}
    \clearpage
}
%===========================================

%===========================================
% INSTRUÇÃO IA: CONFIGURAÇÃO DE CABEÇALHOS (RUNNING HEADS)
% OBJETIVO: Implementar o estilo de cabeçalho "meuestilo" (com reta e títulos).
% IMPLEMENTAÇÃO:
% 1. Estilo de Página (\makepagestyle{meuestilo}): Implementar o código fornecido.
% 2. Geração da Marca (\leftmark): A IA deve redefinir o comando de marcação (\markboth ou \markright) para que \leftmark inclua o rótulo completo ("APÊNDICE A", "ANEXO B", ou "1") para a formatação correta do cabeçalho.
% 3. Páginas de Abertura: A IA deve garantir que a primeira página de cada \chapter use \thispagestyle{plain} ou equivalente, exibindo APENAS o número da página no canto superior direito.
% 4. Aplicação: O estilo de página \pagestyle{meuestilo} será ativado dentro de \mainmatter.
%===========================================
\makepagestyle{meuestilo}
\makeevenhead{meuestilo}{\footnotesize\textit{\thechapter.}\quad\textit{\leftmark}}{}{\footnotesize\thepage}
\makeoddhead {meuestilo}{\footnotesize\textit{\thechapter.}\quad\textit{\leftmark}}{}{\footnotesize\thepage}
\makeheadrule{meuestilo}{\textwidth}{0.3pt}
%\pagestyle{meuestilo} % Será ativado dentro de \mainmatter


\begin{document}

% INSTRUÇÃO IA: Configuração Inicial de Numeração
% 1. Numeração ARÁBICA (1, 2, 3...). NUNCA usar romanos.
% 2. O comando \frontmatter deve ser inserido aqui para ativar a contagem invisível.
%\frontmatter

\folha
\contra

%===========================================
% INSTRUÇÃO IA: ELEMENTOS PRÉ-TEXTUAIS (INÍCIO)
% A IA deve usar o comando modular \capituloCentralizado.
%===========================================

\chapter*{Agradecimento}
\chapter*{Resumo}
\chapter*{Lista de Figuras}
\chapter*{Lista de Tabelas}
\chapter*{Lista de Abreviaturas e Siglas}
\chapter*{Lista de Símbolos}


%===========================================
% INSTRUÇÃO IA: CONFIGURAÇÃO DO SUMÁRIO (CLASSE MEMOIR)
% IMPLEMENTAÇÃO: Centralização do Título, Supressão do Título "SUMÁRIO" no índice, Títulos em CAIXA ALTA, 2 cm de espaçamento horizontal.
%===========================================
\tableofcontents % A IA deve aplicar as configurações acima a este comando.
%===========================================


%===========================================
% INSTRUÇÃO IA: TRANSIÇÃO PARA A PARTE TEXTUAL (CRÍTICO)
% 1. Numeração: ARÁBICA e CONTÍNUA, SEM REINICIAR O CONTADOR.
% 2. Início da Exibição: O comando \mainmatter deve ser inserido ANTES da Introdução. Ele DEVE ativar a exibição VISÍVEL da numeração arábica, o ESTILO DE CABEÇALHO (meuestilo) e o ESTILO DE CAPÍTULO (CapitulosNumerado).
% 3. Introdução (Híbrida): Deve usar o comando modular \capituloCentralizado.
%===========================================
%\mainmatter

% A Introdução deve ser \chapter* (sem número) + \addcontentsline (para Sumário).
\chapter*{Introdução}

% Contagem de Capítulos DEVE continuar sequencialmente (1, 2, 3...), SEM SER REINICIADA após a página \part.
\part{Preparação do Relatório}
\chapter{Desenvolvimento}
\part{Resultado do Relatório}
\chapter{Conclusão}


%===========================================
% INSTRUÇÃO IA: ELEMENTOS PÓS-TEXTUAIS (CRÍTICO)
%
% OBJETIVO: Apêndice/Anexo com numeração alfabética 'A' (A.1, A.2, ...) INDEPENDENTE.
%
% 1. Transição Pós-Textual: O comando \backmatter deve ser inserido antes das Referências.
% 2. Referências: Deve ser um \chapter* (usando \capituloCentralizado).
% 3. Títulos divisórios: Usar \capituloCentralizado, listando no TOC SEM o número da página.
% 4. Numeração: A IA DEVE INSERIR COMANDOS DE RESET DE CONTADOR para garantir que o Anexo comece do 'A'.
% 5. Marca do Cabeçalho: A IA deve garantir que a marca do cabeçalho seja redefinida para incluir a palavra "APÊNDICE" ou "ANEXO".
%===========================================
%\backmatter

\chapter*{Referência}

%\appendix

% A IA deve criar um comando que use \chapter* + \addcontentsline{toc} sem número de página
\chapter*{Apêndices}
\chapter{Título do Apêndice A} % Deve gerar APÊNDICE A
\chapter{Título do Apêndice B} % Deve gerar APÊNDICE B

% A IA deve inserir o comando de RESET do contador AQUI e redefinir a nomenclatura para ANEXO.

% A IA deve criar um comando que use \chapter* + \addcontentsline{toc} sem número de página
\chapter*{Anexos}
\chapter{Título do Anexo C} % Deve gerar ANEXO A
\chapter{Título do Anexo D} % Deve gerar ANEXO B

\end{document} 

%Implemente todos os comandos e configurações comentados, garantindo a total modularidade e o uso exclusivo dos recursos da classe memoir, conforme solicitado. 


%Seria útil estruturar esses modelos diretamente em LaTeX, junto com uma instrução clara para a Inteligência Artificial. É preciso considerar que alguns documentos usam títulos e capítulos estilizados, enquanto outros mantêm o padrão. Isso precisa ser levado em conta desde o início.
%
%Também é necessário diferenciar corretamente as partes pré-textuais, textuais e pós-textuais, pois algumas páginas têm numeração e outras não. Os comandos devem funcionar de forma consistente em todas essas seções, inclusive nas que não possuem número.
%
%O ponto crítico é a estilização de capítulos. Quando se faz muitas adaptações para remover numeração em páginas pré-textuais, personalizar apêndices e anexos, qualquer mudança no estilo do capítulo tende a quebrar o restante do documento. Se o usuário optar por ativar um estilo específico de capítulo, ele deve ser aplicado tanto aos capítulos numerados quanto aos não numerados. E, se o estilo for removido, o documento precisa voltar ao comportamento padrão sem causar conflitos.
%
%O problema atual é que as gambiarras usadas para ligar e desligar numeração acabam afetando outras partes, como sumário, contracapa e folha de rosto. Isso gera inconsistências quando um estilo é ativado ou desativado. Por isso, a estrutura deve ser organizada, modular e fácil de manter. 